\section{Conclusion}\label{conclusion}

\subsection{Summary of Contributions}\label{summary-of-contributions}

This paper presented ADL Lite, an event-first operational ontology for multi-agent concept consensus, and evaluated it through nine experiments establishing correctness properties across chain integrity, status derivation, snapshot round-trip consistency, precondition enforcement, multi-agent auditability, and scalability on real-world financial transaction data. The four principal contributions are restated below with the quantitative results obtained.

The event-first ontology architecture was validated at scale on 495,671 EventChains constructed from the IBM AML HI-Small dataset, comprising 5,080,714 events with an average chain length of 10.3 events and a maximum of 168,508. Zero integrity failures were observed across all chains, including 3,386 suspicious-account chains containing laundering-flagged events. The 100\% derivation accuracy demonstrated in E2 (2,204/2,204 exhaustive event sequences) confirms that status, confidence, and validator identity can be computed reliably from append-only cryptographic chains, eliminating the stored-state consistency bugs endemic to object-first designs \cite{ref3}\cite{ref14}.

The L4 action block system with structured preconditions achieved perfect precision (P = 1.0), perfect recall (R = 1.0), and perfect F1 score (F1 = 1.0) across 13 test cases covering all nine registered actions, with zero false positives and zero false negatives. The \texttt{Comparator} enumeration and \texttt{PreconditionRule} Pydantic model enforce action preconditions through static type checking without runtime \texttt{eval()}, closing a code-injection vector that is present in dynamically evaluated condition systems \cite{ref8}.

The open-source \texttt{adl\_lite} Python toolkit implements the full architecture---parser, validator, \texttt{ActionExecutor}, \texttt{ConsensusEngine}, \texttt{DataImporter}, and unified experiment runner---as a pip-installable package operating over Markdown files in standard Git repositories. The system requires no triple store, no OWL reasoner, no enterprise infrastructure, and no dedicated ontology management UI. Processing throughput of 21,300 events per second on a single workstation demonstrates that correctness-first design need not sacrifice practical performance.

The central takeaway of this work can be stated in one sentence: \emph{by making events the fundamental ontological primitive and deriving all concept state from cryptographically linked append-only chains, ADL Lite achieves deployment-ready lifecycle traceability and multi-agent consensus---under a collaborative-audit threat model---at a cost closer to personal note-taking software than to enterprise knowledge graph platforms.}

\subsection{Future Work}\label{future-work}

Eight directions for future work emerge from the limitations identified in Section 6.4, the trust-model analysis in Section 6.5, and the concurrency discussion in Section 6.1.

\textbf{Real side-effect backends.} The immediate priority is to replace the Lark messaging stub with a production-strength instant-messaging backend and to implement the \texttt{publish} and \texttt{dashboard} side-effect handlers against real services. Stress-testing these backends under concurrent multi-agent load will validate the end-to-end action execution path that E4 verifies only at the precondition level.

\textbf{Complete ActionExecutor--ConsensusEngine integration.} The \texttt{ScriptedHarness} in E5 mutates \texttt{FrontMatter} directly rather than exercising the full event-driven status transition loop. Refactoring the harness to append \texttt{SimEvent} objects to chains and re-derive all front-matter fields will validate the integrated event loop and close the gap between component-level and system-level testing.

\textbf{Full IBM AML dataset expansion.} The current evaluation uses the HI-Small subset (5.08 million events). Scaling to the full 7.6 GB IBM AML dataset will test the architecture at approximately tenfold data volume. This expansion is architectural (scalability stress testing), not domain validation---pattern-detection heuristics remain uncalibrated against expert labels.

\textbf{CRDT-based event chain merging.} The current \texttt{ConsensusEngine} resolves concurrent edits to the same concept through a last-write-wins strategy that selects a single winning event and discards concurrent alternatives. Replacing this with operation-based CRDT semantics---inspired by Automerge---would provide automatic conflict resolution that preserves all concurrent edits as a multi-register rather than choosing a single winner \cite{ref21}. Each concurrent update would be retained as a concurrent branch in the event DAG, with merge functions defined for each payload field (e.g., set-union for \texttt{evidence\_links}, max-wins for \texttt{confidence}). This would eliminate lost updates in high-contention multi-agent scenarios and align the merge semantics with the event-first design principle that no observation should be discarded.

\textbf{Linked Data Proofs for authenticated events.} The most significant trust-model enhancement is cryptographic actor authentication. Each event should be signed using Linked Data Proofs: URDNA2015 canonicalisation of the event payload followed by an ed25519 signature from the actor's private key \cite{ref20}. This binds every event to a non-repudiable identity, closes the impersonation vulnerability identified in Section 6.5, and enables decentralised verification without a central authority. Key management---generation, distribution, rotation, and revocation---must be designed for the multi-agent, multi-tenant context.

\textbf{PROV-O/RDF serialisation and SHACL validation.} Exporting EventChains to W3C PROV-O format \cite{ref19} would enable interoperability with standard provenance tools and repositories. SHACL constraints should be defined for: (i) valid status-transition shapes; (ii) mandatory fields per event type; and (iii) actor-scope binding rules. SHACL validation would complement the imperative \texttt{ActionExecutor} checks with a declarative, standardised constraint language that third-party tools can evaluate independently.

\textbf{Large-scale multi-agent evaluation with conflict metrics.} E5's five-chain pilot must be superseded by systematic simulations measuring: \emph{conflict rate} (frequency of contradictory concept proposals across agents); \emph{time-to-consensus} (distribution of review cycles until validation or deprecation); \emph{fork rate} (proportion of disagreements that result in forked concept variants versus resolved merges); and \emph{provenance completeness} (fraction of status changes captured in the chain versus bypassed via direct mutation). These metrics require hundreds of chains, dozens of simulated agents, and deliberately injected conflicting proposals over extended simulation durations.

\textbf{Domain-level AML evaluation with expert labels.} A proper AML domain evaluation requires expert-annotated laundering labels for a representative sample of the IBM dataset, concept-governance correctness metrics measuring alignment between derived ontology structures and domain-expert judgments, and quantified false-positive and false-negative rates against a validated Financial Action Task Force (FATF) typology. Such an evaluation is a multi-month interdisciplinary effort involving financial-crimes experts and is beyond the scope of a systems-paper contribution, but it is necessary before the AML components could be considered for operational deployment \cite{ref16}.

\textbf{Configurable ontology discovery.} Replacing the simple \texttt{\_id} suffix heuristic with a configurable discovery rule engine---supporting regex patterns, statistical co-occurrence thresholds, and optional LLM-based class name suggestion---will improve the quality of discovered object and link types without requiring the pre-definition overhead of enterprise ontology platforms \cite{ref7}.
